\chapter{Conclusiones y Trabajo Futuro}
\label{cap:conclusiones}

\section{Conclusiones}
\label{sec:conclusiones}

Este trabajo ha validado la viabilidad de la metodología QHRP (Hierarchical Risk Parity cuántico) como una extensión robusta de los métodos de optimización de carteras clásicos. A través de la representación de activos en el espacio de Hilbert, se ha demostrado que la geometría cuántica permite capturar estructuras de riesgo que la correlación lineal ignora.

\subsection{Síntesis de hallazgos}
Los experimentos realizados permiten extraer las siguientes conclusiones principales:

\begin{itemize}
    \item \textbf{Ventaja en rendimiento:} QHRP ha mostrado una capacidad superior de acumulación de retornos ajustados por riesgo en el horizonte fuera de muestra, obteniendo el mejor ratio de Sharpe agregado. Si bien la variabilidad mensual es alta, la estrategia cuántica demuestra una mayor robustez intertemporal.
    
    \item \textbf{Coherencia y robustez estructural:} La segmentación jerárquica resultante no es un artefacto numérico, sino una representación económicamente interpretable. Los tests de permutación confirman la invarianza del método, mientras que el análisis sectorial revela que QHRP agrupa activos por factores de riesgo coherentes (renta variable pura vs. diversificadores).
    
    \item \textbf{Compromiso rendimiento-coste:} La principal barrera para la adopción de QHRP es el coste computacional derivado de la simulación de estados cuánticos. Con una brecha temporal significativa respecto al método clásico, su uso actual se recomienda para análisis estratégico o carteras de tamaño medio (hasta 50 activos) donde la calidad de la segmentación prime sobre la velocidad de ejecución.
\end{itemize}

\subsection{Contribuciones y limitaciones}
La tesis contribuye al estado del arte mediante la reproducción fiel y extensión de la metodología QHRP, aportando validaciones de robustez (permutación) y una primera transferencia a hardware real en dispositivos de IBM. El piloto en hardware, aunque limitado por el ruido de los dispositivos NISQ actuales, confirma que la estructura jerárquica cuántica persiste fuera del entorno de simulación ideal.

No obstante, el estudio reconoce limitaciones intrínsecas: la dependencia de la simulación clásica para el grueso de los experimentos, la ausencia de costes de transacción en el modelo y la sensibilidad del modelo a hiperparámetros como el factor de escala en el mapeo cuántico. Estas limitaciones definen el marco de confianza de los resultados obtenidos.

\section{Trabajo Futuro}
\label{sec:trabajo_futuro}

La investigación abre diversas líneas de desarrollo para consolidar la ventaja cuántica en la gestión de carteras:

\begin{enumerate}
    \item \textbf{Validación sistemática en hardware real:} Superar el piloto preliminar mediante ejecuciones exhaustivas en ordenadores cuánticos reales. Resulta prioritario testear configuraciones de 5 o 6 qubits (correspondientes a las características $P$ del modelo) para evaluar cómo el ruido físico y los errores de puerta afectan a la estabilidad de la métrica de \textit{block contrast} y a la segmentación final.
    
    \item \textbf{Optimización de la escalabilidad:} Investigar la simplificación de los \textit{feature maps} y técnicas de compilación de circuitos que permitan manejar universos de activos más amplios manteniendo la profundidad de circuito en niveles tolerables para el hardware NISQ.
    
    \item \textbf{Métricas de riesgo y adaptabilidad:} Explorar la inclusión de métricas de riesgo no lineales (como el CVaR cuántico) y el desarrollo de versiones dinámicas del algoritmo que se adapten a cambios bruscos de régimen de mercado mediante rebalanceos optimizados.
\end{enumerate}

En conclusión, QHRP se posiciona como una herramienta prometedora que trasciende las limitaciones de la correlación de segundo orden, sentando las bases para una nueva generación de algoritmos de optimización de carteras basados en la ventaja de la representación cuántica.